\subsection{从物理对象到运动方程}

牛顿定律的应用不是把若干力机械地写进公式，而是先说明研究对象、参考系和作用过程。
把一个物体单独取出时，凡是来自系统边界以外的相互作用都要画在受力图上；物体内部各部分
之间的作用不应重复计入。重力由质量和重力场决定，弹簧弹力由形变量决定，支持力、绳张力
和静摩擦力却通常要与运动约束一起求出。因此受力图上可以先标未知量，但不能凭经验预先
认定它们的大小，更不能把“运动方向”误当成“合力方向”。

选定惯性系后，质点的运动方程为
\begin{equation}
 \frac{\dd\vec p}{\dd t}=\sum_i\vec F_i,
 \qquad\text{质量不变时为}\qquad
 m\ddot{\vec r}=\sum_i\vec F_i.
\end{equation}
这是一个矢量方程。若在直角坐标系中使用，就应在每个方向分别投影；若轨迹弯曲，沿切向与
法向分解往往更直接：切向分量改变速率，法向分量改变速度方向。得到加速度以后，还必须结合
初始位置和初速度积分，运动才被唯一确定。由此可见，同一受力方程配上不同初始条件，可以
描述完全不同的运动。

约束把不同物体的运动联系起来。不可伸长轻绳给出绳长不变；对于方向不转动的平面接触，
保持接触意味着物体与平面沿固定法向的加速度相同；纯滚动给出接触点相对速度为零。
曲面接触则不能直接照搬“法向加速度相同”。例如质点沿固定圆轨道运动时，
由 $\vec r\cdot\vec r=R^2$ 先求导得 $\vec r\cdot\vec v=0$，再求导得
$\vec r\cdot\vec a=-v^2$，因此即使轨道静止，质点仍有大小 $v^2/R$ 的向心加速度。
一般约束须先写几何关系，再对时间求导，不能遗漏曲率或转动带来的项。
可靠的列式顺序是：分别画受力图，选取便于投影的方向，
为每个独立物体列动力学方程，再用几何或运动学约束把未知量闭合。方程数量不足时，往往不是
“少列了一条牛顿定律”，而是遗漏了一个约束关系。

牛顿第三定律说明相互作用总是成对出现：$\vec F_{A\to B}=-\vec F_{B\to A}$。这两个力作用在
不同物体上，所以不能在单个物体的受力图中相互抵消；只有把 $A,B$ 作为一个整体研究时，它们
才成为内力并在系统合力中消去。第三定律还要求同时、共线，并与两个物体是否静止无关。
处理多物体问题时，先整体求加速度、再隔离局部求张力或支持力，正是利用了这一结构。
